\parttitle[chemistry]{The Chemistry of Life}{Water, bonds, pH, buffers, and the thermodynamic rules every later chapter obeys.}

\section{Water and the hydrogen bond}

\begin{defbox}[title={DEFINITIONS: bonds and forces}]
\nr{stub-bond-types}{Bond types in biology}
A \tm{covalent bond} shares electrons, $348\ \mathrm{kJ\,mol^{-1}}$ for \ce{C-C}. A
\tm{hydrogen bond} is $20\ \mathrm{kJ\,mol^{-1}}$ and water at $25\degc$ keeps about
$3.4$ of them per molecule.
\end{defbox}

\begin{keybox}[title={MEMORIZE: the water potential equation}]
\nr{stub-water-potential}{Water potential}
\[
\Psi_w=\Psi_s+\Psi_p,\qquad \Psi_s=-iCRT
\]
with $R=8.31\ \mathrm{J\,mol^{-1}\,K^{-1}}$ and $T$ in kelvin. Water moves from high
$\Psi_w$ to low $\Psi_w$; $30\%$ of exam errors here are sign errors.
\end{keybox}

\begin{calcbox}[title={WORKED: buffer pH}]
A buffer holds $0.10\ \mathrm{M}$ \ce{CH3COOH} and $0.25\ \mathrm{M}$ \ce{CH3COO-},
$pK_a=4.76$.
\[
\mathrm{pH}=4.76+\log_{10}\frac{0.25}{0.10}=4.76+0.40=5.16
\]
\ttitle{Check.} More base than acid, so pH must exceed $pK_a$. \tm{$\surd$}
\end{calcbox}

\begin{expbox}[title={EXPERIMENT: Miller and Urey, 1953}]
\nr{stub-miller-urey}{Miller-Urey experiment}
Sparking \ce{CH4}, \ce{NH3}, \ce{H2} and \ce{H2O} for a week yielded glycine, alanine
and $11$ of the $20$ standard amino acids, at up to $2\%$ of the carbon supplied.
\end{expbox}

\begin{warnbox}[title={TRAP: osmolarity is not tonicity}]
\nr{stub-tonicity-trap}{Osmolarity against tonicity}
A solution of $300\ \mathrm{mOsm}$ urea is iso-osmotic to a cell and yet hypotonic,
because urea crosses the membrane. Tonicity counts only the impermeant solutes.
\end{warnbox}

\begin{tipbox}[title={TECHNIQUE: follow the labelled atom}]
Trigger: a question gives a radiolabel. Track the atom, not the molecule.
\begin{itemize}
\item $^{18}$O in water appears in \ce{O2} in photosynthesis.
\item $^{14}$C in \ce{CO2} appears first in 3-phosphoglycerate.
\end{itemize}
\end{tipbox}

\begin{thmbox}[title={PRINCIPLE: the second law in a cell}]
\nr{stub-second-law}{Second law applied to organisms}
An organism decreases its own entropy only by exporting more to its surroundings, so
$\Delta S_{\text{universe}}>0$ always holds. Growth needs $\Delta G<0$ overall.
\end{thmbox}

\begin{obsbox}[title={WORTH NOTICING: the $10\%$ and the $40\%$}]
Trophic transfer is near $10\%$ and respiration captures near $40\%$ of glucose free
energy. Both numbers are efficiencies and both are quoted as ranges in practice.
\end{obsbox}

\subsection{A table that has to shrink}

\begin{tabular}{@{}lp{3.2cm}p{4.4cm}@{}}
\toprule
\textbf{Force} & \textbf{Energy} & \textbf{Where it matters}\\
\midrule
Covalent & $200$--$800\ \mathrm{kJ\,mol^{-1}}$ & Backbones, $50\%$ of all structure\\
Hydrogen & $4$--$30\ \mathrm{kJ\,mol^{-1}}$ & Base pairing, $\alpha$ helices\\
\bottomrule
\end{tabular}
